% R13 THEORY PREWORK -- NEW FILE ONLY; not included by main.tex or supplement.tex.
% Base: branch scgi-ceiling-diagnostic-r1, HEAD bcbb31254c343eb04e8cb96d7c3157bfc9818b36.
% Purpose: paste-ready replacement material for Appendix D.4 after author review.
% This file deliberately does not preserve claims that are false as stated at bcbb312.

\paragraph{Theorem C$'$ (known-carrier calibrated soft-log; exact finite-window bound).}
Fix a frame $n$ and let $\mathcal I_n$ be the set of distinct frames
used in its window.  Let $w_{nk}\ge0$ and
$\sum_{k\in\mathcal I_n}w_{nk}=1$.  Conditional on the gain path
$\ell=(\ell_k)$ and known positive carriers $(b_k)$, assume that the
counts are independent and
\[
 C_k\mid(\ell,b)\sim\operatorname{Pois}\{\lambda_k(\ell_k)\},
 \qquad
 \lambda_k(t)=\Lambda_0e^t b_k+d,
 \qquad d\ge0.
\]
For fixed $\alpha>0$, put
\[
 \psi_\alpha(c)=\log(c+\alpha),\qquad
 m_\alpha(\lambda)
 =\mathbb E_{\operatorname{Pois}(\lambda)}\psi_\alpha(C),
\]
and define
\[
 q_k(t)=m_\alpha\{\lambda_k(t)\},\qquad
 v_k(t)=\operatorname{Var}_{\operatorname{Pois}(\lambda_k(t))}
             \{\psi_\alpha(C)\},
\]
\[
 M_n(t)=\sum_{k\in\mathcal I_n}w_{nk}q_k(t),\qquad
 Y_n=\sum_{k\in\mathcal I_n}w_{nk}\psi_\alpha(C_k).
\]
Let $\Theta_n=[\theta_{n,-},\theta_{n,+}]$ be a compact operating
interval containing every $\ell_k$, $k\in\mathcal I_n$, and suppose
\[
 \underline\kappa_n:=\inf_{t\in\Theta_n}M_n'(t)>0.
\]
If the calibration curve is certified only on
$[\lambda_-,\lambda_+]$ and $d<\lambda_-$, the largest safe interval
before intersection with a physical parameter range is
\[
 \left[
 \max_{k\in\mathcal I_n}
  \log{\lambda_--d\over\Lambda_0b_k},
 \quad
 \min_{k\in\mathcal I_n}
  \log{\lambda_+-d\over\Lambda_0b_k}
 \right].
\]
Nonemptiness is an assumption of the theorem; a clamp cannot repair an
empty calibration-safe interval.
Extend the inverse of $M_n$ to all $y\in\mathbb R$ by the clamped
generalized inverse
\[
 G_n(y)=M_n^{-1}\!\left\{\Pi_{M_n(\Theta_n)}y\right\}.
\]
The exact calibrated estimator is $\widehat\ell_n=G_n(Y_n)$.  Define
\[
 B_n=\sum_{k\in\mathcal I_n}w_{nk}
       \{q_k(\ell_k)-q_k(\ell_n)\},
 \qquad
 V_n=\sum_{k\in\mathcal I_n}w_{nk}^2v_k(\ell_k).
\]
Then, conditionally on $(\ell,b)$,
\begin{equation}
 \boxed{
 \mathbb E\!\left[(\widehat\ell_n-\ell_n)^2\mid\ell,b\right]
 \le {B_n^2+V_n\over\underline\kappa_n^2}.}
 \tag{D.8-R13}
\end{equation}
This is a finite-window statement.  It requires neither stationarity
nor mixing of the known carrier sequence.

\paragraph{Proof.}
The Poisson derivative identity gives
\[
 {d\over d\lambda}\mathbb E f(C)
 =\mathbb E\{f(C+1)-f(C)\}.
\]
Consequently
\[
 q_k'(t)=\Lambda_0e^tb_k\,
   \mathbb E_{\operatorname{Pois}(\lambda_k(t))}
   \log{C+1+\alpha\over C+\alpha}>0,
\]
so $M_n$ is strictly increasing.  Projection onto an interval is
nonexpansive, and the ordinary inverse of $M_n$ is
$\underline\kappa_n^{-1}$-Lipschitz.  Hence its clamped extension is
globally Lipschitz:
\[
 |G_n(y_2)-G_n(y_1)|
 \le {|y_2-y_1|\over\underline\kappa_n}
 \qquad(y_1,y_2\in\mathbb R).
\]
Write
\[
 Y_n-M_n(\ell_n)=B_n+Z_n,
\]
where
\[
 Z_n=\sum_{k\in\mathcal I_n}w_{nk}
 \{\psi_\alpha(C_k)-q_k(\ell_k)\}.
\]
Conditional independence gives
\[
 \mathbb E(Z_n\mid\ell,b)=0,\qquad
 \operatorname{Var}(Z_n\mid\ell,b)=V_n.
\]
Since $G_n\{M_n(\ell_n)\}=\ell_n$, global inverse Lipschitzness yields
\[
 |\widehat\ell_n-\ell_n|
 \le {|B_n+Z_n|\over\underline\kappa_n}.
\]
Taking conditional second moments proves (D.8-R13).
\hfill$\blacksquare$

\paragraph{Clamp diagnostic (not an extra risk term).}
All-zero and out-of-range windows are already covered by
(D.8-R13).  If a separate diagnostic is desired, let
\[
 r_n=\min\{M_n(\ell_n)-M_n(\theta_{n,-}),
           M_n(\theta_{n,+})-M_n(\ell_n)\}.
\]
Then
\[
 \Pr\{Y_n\notin M_n(\Theta_n)\}
 \le\min\left\{1,{B_n^2+V_n\over r_n^2}\right\},
\]
and, when $|B_n|<r_n$,
\[
 \Pr\{Y_n\notin M_n(\Theta_n)\}
 \le\min\left\{1,{V_n\over(r_n-|B_n|)^2}\right\}.
\]
No term $D_\Theta^2\Pr(\mathcal E_n^c)$ is needed in the risk bound.

\paragraph{Approximate calibration and numerical inversion.}
Let $\widetilde M_n$ be a continuous strictly increasing approximation
on the same interval
$\Theta_n$, with
\[
 \sup_{t\in\Theta_n}|\widetilde M_n(t)-M_n(t)|
 \le\varepsilon_{\mathrm{cal},n}.
\]
Let $\widetilde G_n$ be its clamped generalized inverse, and suppose
the numerical solver returns $\widehat\ell_n^{\mathrm{num}}$ satisfying
\[
 |\widehat\ell_n^{\mathrm{num}}-\widetilde G_n(Y_n)|
 \le\varepsilon_{\mathrm{bis},n},
\]
where $\varepsilon_{\mathrm{bis},n}$ is measured in the log-parameter
domain.  A conservative implementation bound is
\begin{equation}
 \boxed{
 \left\{\mathbb E[(\widehat\ell_n^{\mathrm{num}}-\ell_n)^2
                    \mid\ell,b]\right\}^{1/2}
 \le {\sqrt{B_n^2+V_n}\over\underline\kappa_n}
 +{\varepsilon_{\mathrm{cal},n}\over\underline\kappa_n}
 +\varepsilon_{\mathrm{bis},n}.}
 \tag{D.8b-R13}
\end{equation}
Indeed, using the same endpoint-clamped inverse convention for both
curves, inverse bracketing gives
\[
 G_n(y-\varepsilon_{\mathrm{cal},n})
 \le\widetilde G_n(y)
 \le G_n(y+\varepsilon_{\mathrm{cal},n}).
\]
If tabulation and interpolation have separate certified forward-error
budgets, replace $\varepsilon_{\mathrm{cal},n}$ by their sum.  The
quantities that have the same units are
\[
 \varepsilon_{\mathrm{cal},n}/\underline\kappa_n
 \quad\hbox{and}\quad \varepsilon_{\mathrm{bis},n};
\]
one must not add them in the response domain and divide both by
$\underline\kappa_n$.

\paragraph{A theorem-level interpolation certificate.}
The reported midpoint probe is an empirical diagnostic, not by itself
a uniform bound on every point of the calibration interval.  Put
\[
 x=\log\lambda,\qquad f_\alpha(x)=m_\alpha(e^x).
\]
For a log-intensity grid of maximum spacing $\Delta_x$, piecewise-linear
interpolation from exact node values obeys
\begin{equation}
 \|f_\alpha-I_{\Delta_x}f_\alpha\|_\infty
 \le {\Delta_x^2\over8}\|f_\alpha''\|_\infty. \tag{D.cal-R13}
\end{equation}
With $\Delta g(c)=g(c+1)-g(c)$, the Poisson difference identities give
\[
 f_\alpha''(x)
 =\lambda\,\mathbb E\Delta\psi_\alpha(C)
 +\lambda^2\,\mathbb E\Delta^2\psi_\alpha(C),
 \qquad \lambda=e^x.
\]
Hence a certified bound
\[
 K_{\alpha,2}\ge
 \sup_{\lambda\in[\lambda_-,\lambda_+]}
 \left|\lambda\mathbb E\Delta\psi_\alpha(C)
 +\lambda^2\mathbb E\Delta^2\psi_\alpha(C)\right|
\]
and a certified node-evaluation error $\varepsilon_{\mathrm{tab}}$
yield
\[
 \varepsilon_{\mathrm{cal}}
 \le\varepsilon_{\mathrm{tab}}
 +{K_{\alpha,2}\Delta_x^2\over8}.
\]
The expectations and Poisson tails may be enclosed analytically or by
interval arithmetic.  Until such an enclosure is generated, the measured
$1.6410165\times10^{-6}$ midpoint maximum must be labeled empirical; the
measured log-parameter bisection bound is
$1.5396755\times10^{-17}$.
For $J$ midpoint bisections on a log-parameter interval of diameter
$D_n$, the parameter-domain numerical error is certified by
\[
 \varepsilon_{\mathrm{bis},n}\le {D_n\over2^{J+1}}.
\]

\paragraph{H\"older bias and truncated windows.}
Suppose $\ell_k=f(k/N)$ and, for $0<\beta\le1$,
\[
 |f(x)-f(y)|\le L|x-y|^\beta.
\]
Let
\[
 h_n=\max_{k\in\mathcal I_n}{|k-n|\over N},\qquad
 \overline\kappa_k=\sup_{t\in\Theta_n}q_k'(t),\qquad
 W_{\mathrm{eff},n}={1\over\sum_kw_{nk}^2}.
\]
Put $\kappa_{\max,n}=\max_{k\in\mathcal I_n}\overline\kappa_k$.
Then
\[
 |B_n|
 \le L\sum_{k\in\mathcal I_n}w_{nk}\overline\kappa_k
       \left({|k-n|\over N}\right)^\beta
 \le\kappa_{\max,n}Lh_n^\beta,
\]
and, if $v_k(\ell_k)\le\bar\sigma_\alpha^2$,
\[
 V_n\le{\bar\sigma_\alpha^2\over W_{\mathrm{eff},n}}.
\]
Therefore
\begin{equation}
 \mathbb E(\widehat\ell_n-\ell_n)^2
 \le {1\over\underline\kappa_n^2}
 \left\{\kappa_{\max,n}^2L^2h_n^{2\beta}
       +{\bar\sigma_\alpha^2\over W_{\mathrm{eff},n}}\right\}.
 \tag{D.9-R13}
\end{equation}
For a truncated-renormalized equal-weight window,
$W_{\mathrm{eff},n}=|\mathcal I_n|$.  Thus an edge window has a larger
variance constant than an interior window, but the same rate.

\paragraph{Low-photon expansion: the general leading term.}
Let $d=0$ and define
\[
 \lambda_{k,n}^{0}=\Lambda_0e^{\ell_n}b_k,\qquad
 \lambda_k^\star=\Lambda_0e^{\ell_k}b_k,
\]
\[
 A_n=\sum_kw_{nk}\lambda_{k,n}^{0},\qquad
 D_n=\sum_kw_{nk}^2\lambda_k^\star.
\]
Fix $\alpha>0$, put $c_\alpha=\log(1+1/\alpha)$, and suppose
\[
 \max_{k\in\mathcal I_n}\{\lambda_{k,n}^{0},\lambda_k^\star\}
 \le\varepsilon.
\]
The fixed-$\alpha$ Poisson expansions give, uniformly over arbitrary
positive carrier heterogeneity,
\[
 M_n'(\ell_n)=c_\alpha A_n\{1+O_\alpha(\varepsilon)\},
\]
\[
 V_n=c_\alpha^2D_n\{1+O_\alpha(\varepsilon)\},
\]
and therefore
\begin{equation}
 \boxed{
 {V_n\over M_n'(\ell_n)^2}
 ={D_n\over A_n^2}\{1+O_\alpha(\varepsilon)\}.}
 \tag{D.10-R13}
\end{equation}
For $m=|\mathcal I_n|$ equal-weight distinct frames, write
\[
 I_{n,\mathrm{ref}}=\sum_k\lambda_{k,n}^{0},\qquad
 I_{n,\mathrm{act}}=\sum_k\lambda_k^\star.
\]
Then
\[
 {V_n\over M_n'(\ell_n)^2}
 ={I_{n,\mathrm{act}}\over I_{n,\mathrm{ref}}^2}
 \{1+O_\alpha(\varepsilon)\}.
\]
If $h_{\ell,n}=\max_k|\ell_k-\ell_n|$, then
\[
 e^{-h_{\ell,n}}
 \le {I_{n,\mathrm{act}}\over I_{n,\mathrm{ref}}}
 \le e^{h_{\ell,n}}.
\]
Only under local gain flatness $h_{\ell,n}=o(1)$ may the leading term
be simplified to
\[
 {1+O_\alpha(\varepsilon+h_{\ell,n})
  \over I_{n,\mathrm{ref}}}
 ={1+O_\alpha(\varepsilon+h_{\ell,n})
  \over I_{n,\mathrm{act}}}.
\]
When the gain is constant in the window,
$I_{n,\mathrm{act}}=I_{n,\mathrm{ref}}=m\bar\lambda_n$, and only then
does this become $\{1+O_\alpha(\varepsilon)\}/(m\bar\lambda_n)$.
Carrier comparability is not needed; local gain flatness is.

\paragraph{Counterexample to the unqualified $1/(W\bar\lambda)$ claim.}
Take two equal-weight frames with equal carriers, target $\ell_n=0$,
the other gain $\ell=\log2$, and reference intensity $\varepsilon$.
Then
\[
 I_{n,\mathrm{ref}}=2\varepsilon,\qquad
 I_{n,\mathrm{act}}=3\varepsilon,
\]
so (D.10-R13) gives $3/(4\varepsilon)$, whereas the unqualified formula
gives $1/(2\varepsilon)$.  Their ratio is $3/2$ and does not converge
to one as $\varepsilon\downarrow0$.

\paragraph{From a variance factor to estimator asymptotics.}
Equation (D.10-R13) is an algebraic expansion of the local
variance-to-sensitivity factor.  To identify it with the estimator's
leading MSE, one additionally needs a local variance regime such as
\[
 I_{n,\mathrm{act}}\to\infty,\qquad
 |B_n|=o(\sqrt{V_n}),
\]
a Lindeberg condition, negligible clamp probability, and inverse-map
linearization, for example
\[
 {\sup_{\Theta_n}|M_n''|\sqrt{V_n}
  \over M_n'(\ell_n)^2}\to0.
\]
For bounded total information the compact clamp makes risk saturate;
one cannot globally identify the MSE with $1/I_n$.

\paragraph{First-moment obstruction for $1<\beta\le2$.}
Let $u_k=(k-n)/N$ and suppose locally
\[
 \ell_k-\ell_n=a_nu_k+r_{nk},\qquad
 |r_{nk}|\le L|u_k|^\beta.
\]
If $Q_j$ uniformly bounds the first two derivatives of $q_k$, Taylor's
theorem gives
\begin{align}
 |B_n|\le{}&
 |a_n|\left|\sum_kw_{nk}q_k'(\ell_n)u_k\right|
 +Q_1L\sum_kw_{nk}|u_k|^\beta\notag\\
 &+{Q_2\over2}\sum_kw_{nk}
   \{|a_n||u_k|+L|u_k|^\beta\}^2.
 \tag{D.11-R13}
\end{align}
Thus a local-constant estimator has $O(h_n^\beta)$ deterministic bias
only if the slope-weighted carrier first moment satisfies
\[
 |a_n|\left|\sum_kw_{nk}q_k'(\ell_n)u_k\right|
 \lesssim\underline\kappa_n h_n^\beta.
\]
An internal symmetric window with a homogeneous carrier satisfies this
condition.  A heterogeneous carrier generally does not.  A truncated
one-sided edge window fails even with a homogeneous carrier: for the
first current window $k=0,\ldots,32$, target $n=0$,
\[
 {1\over33}\sum_{k=0}^{32}{k\over N}={16\over N},
\]
which is first order in the window radius.  Therefore the current
local-constant estimator cannot support an all-frame uniform
$\beta>1$ rate.  A valid claim must be restricted to internal balanced
windows or must introduce a separately proved local-polynomial edge
estimator.  For $\beta>2$, moments through order $\lfloor\beta\rfloor$
must be cancelled.

\paragraph{Scope: oracle carriers, known carrier law, and blind recovery.}
The theorem above is conditional on the realized per-frame carriers
$(b_k)$ being known.  This is precisely the oracle-known-carrier arm
implemented in the low-photon experiment.  It does not establish blind
low-photon recovery from an unknown carrier law.

If an unobserved carrier $B$ has a fully known law $F_B$, one may instead
calibrate the different marginal curve
\[
 \bar q_{F_B}(t)=\mathbb E_{B\sim F_B}
 m_\alpha(\Lambda_0e^tB+d).
\]
That estimator needs a separate variance or long-run-variance theorem;
it is not the per-frame known-carrier estimator above.  If $F_B$ is
unknown, $\bar q_{F_B}$ is an unknown nonlinear function.  Unlike the
high-count logarithmic model, it cannot generally be absorbed into one
additive gauge constant.  For example, with $s=\Lambda_0e^t$,
\[
 m_\alpha(sB)=\log\alpha+c_\alpha sB
 +a_{2,\alpha}s^2B^2+O(s^3B^3),
\]
where
\[
 a_{2,\alpha}={1\over2}
 \log{\alpha(\alpha+2)\over(\alpha+1)^2}\ne0.
\]
The laws $B\equiv1$ and
$\Pr(B=1/2)=\Pr(B=3/2)=1/2$ have the same mean but different second
moments, hence different soft-log response curves.  This proves that
the known-carrier theorem and one homogeneous calibration inverse do
not automatically yield an unknown-law blind theorem.  It does not
rule out every possible joint semiparametric method.

\paragraph{Centered-log quotient risk.}
Let $e=\widehat\ell-\ell$ and
$P_N=I-N^{-1}\mathbf1\mathbf1^\top$.  Since $P_N$ is an orthogonal
projection,
\[
 {1\over N}\mathbb E\|P_Ne\|_2^2
 \le {1\over N}\sum_{n=1}^N\mathbb Ee_n^2.
\]
Thus the pointwise bounds transfer directly to record-centered
log-gain loss; overlapping windows need not be independent for this
expectation bound.

% End of paste-ready D.4 replacement material.
